% ======================================================================
% 第 13 章：临界指数与标量场论
% （原书第 439-468 页；PDF 第 462-491 页）
% ======================================================================
\chapter{临界指数与标量场论}

跑动耦合常数与重整化群流的思想为我们提供了一种新的语言，用以讨论标量场论的定性行为。在我们对 $\phi^4$ 理论的首次讨论中，耦合常数的每一个取值——更一般地，势能的每一种形式以及每一个时空维度——都对应着一个需要分别探究的问题。但在第~12 章中我们看到，具有不同耦合常数值的 $\phi^4$ 理论之间由重整化群流联系起来，而这些流的图样随时空维度连续变化。在此背景下，提出如下非常一般的问题是有意义的：$\phi^4$ 理论如何随维度变化？本章将详细回答这一问题。

我们分析的核心要素将是第~12.5 节讨论的威尔逊-费舍尔不动点。该不动点存在于 $d<4$ 的时空维度中；在这些维度中它控制着无质量 $\phi^4$ 理论的重整化群流。标量场论根据质量参数 $m^2$ 的符号而具有明显对称性或自发破缺的对称性。在 $m^2=0$ 附近，理论表现出具有反常维度的标度行为，其数值由重整化群方程决定。对于 $d>4$，威尔逊-费舍尔不动点消失，只剩下自由场不动点。此时理论同样表现出两个不同的相，但相变处的行为现在由自由场不动点附近的重整化群流决定，因此标度律正是那些由简单量纲分析得出的标度律。

将这些结果延拓到欧几里得空间，对磁体与流体中的相变理论具有重要意义。如我们在前一章所讨论的，重整化群的思想意味着相变点附近热力学量的幂律行为由欧几里得 $\phi^4$ 理论中的关联函数行为所决定。上一段的结论进而对临界标度律给出如下推论：对于处于维度 $d>4$ 的空间中的统计系统，标度律正是那些由简单量纲分析得出的标度律。这些预言正是我们在第~8 章讨论过的 Landau 理论的预言。另一方面，对于 $d<4$，临界标度律发生修正，其修正方式可以用重整化群加以计算。

在 $d=4$ 时，我们处于两种标度行为之间的边界上。这正对应 $\phi^4$ 理论恰好可重整的情形。此时量纲分析的预言会得到修正，但仅以对数的方式修正。我们将在第~13.2 节具体分析这一情形。

虽然并不明显，$d=2$ 的情形提供了另一个边界。在这里，向自发对称性破缺的转变由另一个量子场论描述，该理论在二维中变为可重整的。在第~13.3 节中，我们将介绍这一称为\emph{非线性西格马模型}的理论，并展示其重整化群行为如何与 $\phi^4$ 理论的光滑地融合。通过综合本章的所有结果，我们将在整个时空维度范围内定量地理解 $\phi^4$ 理论以及临界现象的行为。

\section{临界指数理论}

在第~12 章末尾，我们利用标量场论重整化群的性质，对热力学系统临界点附近的关联行为作出了预言。我们论证了，当趋近临界点时，关联的范围——关联长度 $\xi$——应按标度律 \eqref{eq:12.141} 增大至无穷。该式中的指数称为 $\nu$，应仅依赖于序参量的对称性。我们还进一步论证，该指数与 $\phi^4$ 理论中某个局域算符的反常维度有关，并且可以由 Feynman 图计算得出。在本节中，我们将证明类似的结论更一般地适用于与临界点相关的大量标度律。

首先，我们将系统地定义一组\emph{临界指数}，即描述临界点附近热力学行为的标度律的指数。然后我们将利用 Callan-Symanzik 方程证明，所有这些指数都可以归结为两个基本的反常维度。最后，我们将把量子场论的这一非凡预言与实验进行比较。

在提出一组临界标度律时，我们从序化场涨落的关联函数行为开始。为具体起见，我们将采用适用于磁体的语言，如同第~8 章那样。我们将把经典热力学期望值当作欧几里得量子场论中的关联函数来计算，如第~9.3 节所解释。涨落场将被称为自旋场 $s(\mathbf{x})$，其积分是磁化强度 $M$，而与 $s(\mathbf{x})$ 耦合的外场将被称为磁场 $H$。（考虑到磁化强度，本节中我们将 Callan-Symanzik 方程中的重整化标度记作 $\mu$。）

定义两点关联函数为
\begin{equation}
  G(x) = \langle s(x)s(0)\rangle ,
  \label{eq:13.1}
\end{equation}
或者，如果我们处于 $\langle s(x)\rangle \neq 0$ 的磁化相中，则以连通期望值定义。远离临界点时，$G(x)$ 应按如下方式指数衰减：
\begin{equation}
  G(x) \sim \exp[-|x|/\xi] .
  \label{eq:13.2}
\end{equation}
趋近临界点由参数
\begin{equation}
  t = \frac{T - T_c}{T_c} .
  \label{eq:13.3}
\end{equation}
刻画。于是我们预期，当 $t\to 0$ 时关联长度应增大至无穷。按如下公式定义指数 $\nu$，\eqref{eq:12.141}：
\begin{equation}
  \xi \sim |t|^{-\nu} .
  \label{eq:13.4}
\end{equation}
恰在 $t=0$ 处，关联函数应仅以幂律衰减。按如下公式定义指数 $\eta$：
\begin{equation}
  G(x) \sim \frac{1}{|x|^{d-2+\eta}} ,
  \label{eq:13.5}
\end{equation}
其中 $d$ 是欧几里得空间的维度。

热力学量在临界点附近的行为定义了许多额外的指数。典型地，热力学系统的比热随 $t\to 0$ 而发散；通过固定外场 $H=0$ 时的比热公式定义指数 $\alpha$：
\begin{equation}
  C_H \sim |t|^{-\alpha} .
  \label{eq:13.6}
\end{equation}
由于序在 $t=0$ 处建立，零场下的磁化强度从负方向趋近 $t\to 0$ 时趋于零。按如下公式定义指数 $\beta$（不要与 Callan-Symanzik 函数混淆）：
\begin{equation}
  M \sim |t|^{\beta} .
  \label{eq:13.7}
\end{equation}
即使在 $t=0$ 处，在非零磁场下磁化强度也不为零。将磁化强度随 $H\to 0$ 趋于零的规律写成关系式
\begin{equation}
  M \sim H^{1/\delta} .
  \label{eq:13.8}
\end{equation}
最后，磁化率在临界点发散；我们将这一发散写成关系式
\begin{equation}
  \chi \sim |t|^{-\gamma} .
  \label{eq:13.9}
\end{equation}
式 \eqref{eq:13.4}--\eqref{eq:13.9} 定义了一组临界指数 $\alpha$、$\beta$、$\gamma$、$\delta$、$\eta$、$\nu$，它们可以在一系列热力学系统中由实验测量。%
\footnote{更多临界指数及其关系参见 M.~E.~Fisher, \emph{Repts.\ Prog.\ Phys.} \textbf{30}, 615 (1967)。}

在第~12 章中，我们遵循 Wilson 的观点论证了，处于临界点附近的热力学系统可以用欧几里得量子场论描述。在原子尺度上，该量子场论的拉格朗日量可能很复杂；然而，当我们把短尺度自由度积分掉之后，这一拉格朗日量会简化。如果调整理论的某个参数以确保长程关联的存在，拉格朗日量就必须紧密趋近重整化群的一个不动点。一般地，拉格朗日量将沿一个不稳定的（相关的）方向趋近不动点，该方向对应 $\phi^4$ 理论的质量参数。在 $d<4$ 时，这就是威尔逊-费舍尔不动点。在 $d>4$ 时，它是自由场不动点。为具体起见，在以下的讨论中我们将假设 $d<4$。

\subsection{自旋关联函数的指数}

在这一框架下，我们可以研究自旋-自旋关联函数 $G(x)$ 的行为。由刚才回顾的论证，$G(x)$ 与欧几里得标量场论的两点关联函数成正比。前一章引入的技术可以直接应用。关联函数满足 Callan-Symanzik 方程，即式 (12.125)，
\begin{equation}
  \Bigl[ \mu\frac{\partial}{\partial\mu} + \sum_i \beta_i\frac{\partial}
    {\partial p_i} + \gamma \Bigr]\ G(x; \mu, \{p_i\}) = 0 .
  \label{eq:13.10}
\end{equation}
这里我们把 $\phi^4$ 耦合 $\lambda$ 也归入广义耦合 $p_i$ 之中。

由量纲分析，在 $d$ 维中，
\begin{equation}
  G(x) = \frac{1}{|x|^{d-2}}\ g\bigl( |x|\mu, \{p_i\} \bigr) ,
  \label{eq:13.11}
\end{equation}
其中 $g$ 是无量纲参数的任意函数。（这就是下述陈述的傅里叶变换：$G(p) \sim p^{-2}$ 乘以一个无量纲函数。）以此为出发点，我们可以用第~12.3 节的方法求解 Callan-Symanzik 方程 \eqref{eq:13.10}，得到
\begin{equation}
  G(x) = \frac{1}{|x|^{d-2}}\ h\bigl( \{p_i(x)\} \bigr)\ \exp\Bigl[
    \int_{1/|x|}^{\mu} d\log(|x'|)\ \gamma\bigl( \{p(x')\} \bigr) \Bigr] ,
  \label{eq:13.12}
\end{equation}
其中 $h$ 是无量纲的初始条件。跑动耦合常数 $p_i$ 满足微分方程
\begin{equation}
  \frac{d}{d\log(1/\mu|x|)}\ p_i = \beta_i\bigl( \{p_j\} \bigr) .
  \label{eq:13.13}
\end{equation}
我们在第~12.5 节研究过该方程的解。我们在那里看到，对于趋近威尔逊-费舍尔不动点附近的流，质量算符的无量纲系数在向大距离移动时增大，而其他无量纲参数变小。设 $\lambda^{*}$ 为不动点的位置。于是我们可以更明确地写出
\begin{equation}
  p_m = p_m(\mu|x|)^{-2+\gamma_{\phi^2}(\lambda^{*})},
  \qquad p_i = p_i(\mu|x|)^{-\Delta_i} ,
  \label{eq:13.14}
\end{equation}
其中对 $i \neq m$ 有 $\Delta_i > 0$。如果把 $\lambda$ 对不动点的偏离也当作某个 $p_i$，通过定义
\begin{equation}
  p_\lambda = \lambda - \lambda^{*} ,
  \label{eq:13.15}
\end{equation}
该参数的重要性也随 $|x|$ 的幂次而下降，正如我们在式 (12.96) 中所证明的。用第~12.1 节的语言来说，除 $p_m$ 外，所有参数 $p_i$ 都乘以无关算符，而 $p_m$ 乘以一个相关算符。

为趋近临界点，我们调整底层理论的参数，使得在接近原子尺度的某个标度 $(1/\mu)$ 上 $p_m \ll 1$。如果 $p_m$ 是通过调节热力学系统的温度来调整的，那么 $p_m \sim t$。只要存在一个距离尺度区域，在该区域中 $p_m$ 保持很小而其余 $p_i$ 可以忽略，临界标度律就有效。于是可以这样计算标度律：以 \eqref{eq:13.14} 给出的 $p_m$、并把其余 $p_i$ 设为零，来评估 Callan-Symanzik 方程的解。可以证明，对这一近似的修正正比于 $t$ 的正幂次。

在这一近似中，我们应在 $p_\lambda = 0$ 处，即不动点处，评估 \eqref{eq:13.12} 中的函数 $\gamma(\lambda)$。利用该值以及 $p_m$ 的解，式~\eqref{eq:13.12} 变为
\begin{equation}
  G(x) = \frac{1}{|x|^{d-2+2\gamma(\lambda^{*})}}\ h\bigl(
    p_m|x|^{2-\gamma_{\phi^2}(\lambda^{*})} \bigr) .
  \label{eq:13.16}
\end{equation}
该方程蕴含标度律 \eqref{eq:13.5} 和 \eqref{eq:13.4}：当 $h$ 的宗量足够小时，$G(x)$ 满足式~\eqref{eq:13.5}，其中
\begin{equation}
  \eta = 2\gamma(\lambda^{*}) .
  \label{eq:13.17}
\end{equation}
在大距离处，$h$ 必须指数衰减，因为这个函数来源于标量场传播子。由 $h$ 的宗量，我们推知该指数必然具有如下形式
\begin{equation}
  \exp\Bigl[ -|x|\bigl( p_m \bigr)^{\nu} \Bigr] ,
  \label{eq:13.18}
\end{equation}
其中，如同 \eqref{eq:12.139}，
\begin{equation}
  \nu^{-1} = 2 - \gamma_{\phi^2}(\lambda^{*}) .
  \label{eq:13.19}
\end{equation}
这恰是标度律 \eqref{eq:13.2}、\eqref{eq:13.4}，其中 $\nu$ 由算符 $\phi^2$ 的反常维度给出。

\subsection{热力学函数的指数}

热力学临界标度律可以用类似的方式推导，即研究宏观热力学变量的标度行为。这些量由 Gibbs 自由能导出，或者用量子场的语言说，由标量场论的有效势导出。由于有效势——更一般地，有效作用量——由关联函数构造，这些量的标度行为可以从关联函数的标度律推导出来。

Gibbs 自由能 $G(M, t)$ 与有效势 $V_{eff}(\phi_{cl})$ 通过将经典场 $\phi_{cl}$ 等同为磁化强度 $M$、并将外场 $H$ 等同为 $G$ 对 $M$ 的导数而联系起来。$G$ 的标度行为来源于这样一个事实：它是关联函数的生成泛函。在标度区域中，得到标度公式
\begin{equation}
  G(M, t) = M^{1+\delta}\ h\bigl( t\,M^{-1/\beta} \bigr) ,
  \label{eq:13.35}
\end{equation}
其中当 $t\to 0$ 时 $h$ 具有光滑极限。表示该公式的一种等价方式是
\begin{equation}
  G(M, t) = t^{\beta(1+\delta)}\ f\bigl( M\,t^{-\beta} \bigr) .
  \label{eq:13.36}
\end{equation}
热力学量的标度律立即由这些关系得出。沿 $t=0$ 的直线，由 \eqref{eq:13.35} 得到
\begin{equation}
  H = \frac{\partial G}{\partial M} = h(0)\,M^{\delta} ,
  \label{eq:13.37}
\end{equation}
这恰是 \eqref{eq:13.8}。在临界温度以下，我们通过对 $M$ 极小化 $G$ 来得到磁化强度的非零值。在标度区域中，这一极小出现在 \eqref{eq:13.36} 中函数 $f(m)$ 的极小值 $m_0$ 处。这给出关系 \eqref{eq:13.7}，其形式为
\begin{equation}
  M\,t^{-\beta} = m_0 .
  \label{eq:13.38}
\end{equation}
如果在 $T_c$ 以上且零场下工作，$f$ 的极小必然出现在 $M=0$ 处。于是
\begin{equation}
  G(t) \sim t^{\beta(1+\delta)} .
  \label{eq:13.39}
\end{equation}
为计算比热，我们对温度求两次导数；这给出标度律 \eqref{eq:13.6}，其中
\begin{equation}
  2 - \alpha = \beta(1+\delta) = \frac{d}{2 - \gamma_{\phi^2}(\lambda^{*})} .
  \label{eq:13.40}
\end{equation}
最后，我们必须构造磁化率的标度律。由 \eqref{eq:13.36}，非零 $t$ 时 $H$ 的标度律为
\begin{equation}
  H = \frac{\partial G}{\partial M} = t^{\beta\delta}\ f'\bigl( M\,t^{-\beta}
    \bigr) .
  \label{eq:13.41}
\end{equation}
该关系的逆是标度律
\begin{equation}
  M = t^{\beta}\ c\bigl( H\,t^{-\beta\delta} \bigr) .
  \label{eq:13.42}
\end{equation}
于是零场下的磁化率为
\begin{equation}
  \chi = \Bigl( \frac{\partial M}{\partial H} \Bigr)_t
  = c'(0)\,t^{-\beta(\delta-1)} .
  \label{eq:13.43}
\end{equation}
因此，我们确认式~\eqref{eq:13.9}，其中
\begin{equation}
  \gamma = (\delta-1)\beta = 2\bigl( 1 - \gamma_{\phi^2}(\lambda^{*}) \bigr) .
  \label{eq:13.44}
\end{equation}

我们现在已经用 Callan-Symanzik 函数给出了所有各种临界指数的显式表达式。当维度 $d$ 从下方趋近 4 时，不动点 $\lambda^{*}$ 趋于零。于是六个临界指数趋于它们在简单量纲分析中会取的值：
\begin{equation}
  \eta = 0;\qquad \nu = \tfrac{1}{2};\qquad \alpha = 0;\qquad
  \beta = \tfrac{1}{2};\qquad \gamma = 1;\qquad \delta = 3 .
  \label{eq:13.45}
\end{equation}
\eqref{eq:13.45} 中给出的 $\eta$、$\nu$ 和 $\beta$ 的值正是我们在第~8 章中从临界现象的 Landau 理论推出的值，这并不令人惊讶。其余的值同样可以证明源于 Landau 理论。重整化群分析告诉我们如何系统地修正 Landau 理论的预言，以恰当地计入自旋场的大尺度涨落。

注意，与热力学量相关的所有指数，都是由与关联函数相关的指数相同的要素构造的。从场论的观点看，这是显然的，因为场论中所有的标度律最终都必须来自算符 $\phi^2(x)$ 和 $\phi^4(x)$ 的反常维度，它们正是 $\gamma_{\phi^2}(\lambda^{*})$ 和 $\gamma(\lambda^{*})$。然而，这一结果有一个有趣的实验推论：它蕴含临界指数之间的\emph{与模型无关的关系}。例如，在任何一个具有临界点的系统中，该理论预言
\begin{equation}
  \alpha = 2 - d\nu,\qquad \beta = \tfrac{1}{2}(d-2+\eta)\nu .
  \label{eq:13.46}
\end{equation}
这些关系检验了把临界点等同于重整化群流不动点的这一总体框架。

此外，临界现象的场论方法预言临界指数是\emph{普适的}，其含义是：在 $T\to T_c$ 极限下趋近同一个标量场不动点的凝聚态系统中，这些指数取相同的值。

\subsection{临界指数的数值}

最后，标量场论实际上预言了 $\gamma(\lambda^{*})$ 和 $\gamma_{\phi^2}(\lambda^{*})$ 的值，这既可以通过第~12.5 节所述的按 $\epsilon = 4-d$ 幂次的展开，也可以通过把 $\beta$ 和 $\gamma$ 函数按 $\lambda$ 的幂次直接展开得到。我们可以利用这些表达式来对临界指数作出定量的预言。在第~12.5 节末尾，我们曾在式~\eqref{eq:12.143} 中给出 $\nu$ 展开的前两项，那是此类预言的一个例子。我们现在回到这一问题，给出所有临界指数的场论预言。

在第~12.5 节末尾的讨论中，我们指出，具有不同涨落自旋分量数的磁体，其临界指数观测值不同。一个乐观的假设是：任何具有 $N$ 个涨落自旋分量——更一般地说，在临界点具有 $N$ 个涨落热力学变量——的热力学系统，都将由具有 $N$ 个标量场的同一不动点场论所描述。这一不动点的自然候选是第~11 章讨论的 $O(N)$ 对称 $\phi^4$ 理论的威尔逊-费舍尔不动点。我们现在将描述在该理论中对一般 $N$ 值的临界指数计算。

作为第一步，我们应该计算四个维度中函数 $\beta(\lambda)$、$\gamma(\lambda)$ 和 $\gamma_{\phi^2}(\lambda)$ 的值。这一计算与第~12 章中对普通 $\phi^4$ 理论所做的分析平行，因此我们将只指出本例需要作出的改动。与普通 $\phi^4$ 理论一样，无质量 $O(N)$ 对称理论的传播子在单圈阶不接收场强修正，因此 $\gamma(\lambda)$ 中的单圈项再次为零。在习题~13.2 中，我们计算 $O(N)$ 对称 $\phi^4$ 理论中对 $\gamma(\lambda)$ 的首个、双圈的贡献：
\begin{equation}
  \gamma = (N+2)\ \frac{\lambda^2}{(4\pi)^4}\ \frac{1}{6}
    + \mathcal{O}(\lambda^3) .
  \label{eq:13.47}
\end{equation}
$\phi^4$ 理论中 $\beta$ 函数的单圈贡献来自单圈顶点抵消项 $\delta_\lambda$，如式 (12.44) 所给。对于 $O(N)$ 对称情形，我们在第~11.2 节计算了相应顶点抵消项的发散部分；由式~\eqref{eq:11.22}，
\begin{equation}
  \delta_\lambda = \lambda^2(N+8)\ \frac{\Gamma(2-d/2)}{(4\pi)^2}
    + \cdots .
  \label{eq:13.48}
\end{equation}
按照通向式~\eqref{eq:12.46} 的逻辑，或利用式~\eqref{eq:12.54}，我们得到
\begin{equation}
  \beta = (N+8)\ \frac{\lambda^2}{16\pi^2} + \mathcal{O}(\lambda^3) .
  \label{eq:13.49}
\end{equation}
如果我们令 $N = 1$ 并作替换 $\lambda\to\lambda/6$，如式~\eqref{eq:11.5} 下方所指出的，这就约化为 $\phi^4$ 理论的 $\beta$ 函数。最后，为计算 $\gamma_{\phi^2}$，我们必须重复第~12.4 节末尾所做的计算。如果考虑（代替 \eqref{eq:12.113}）格林函数 $\langle\phi^i(p)\phi^j(q)\phi^2(k)\rangle$，并把 $\phi^4$ 理论的顶点换成由拉格朗日量 \eqref{eq:11.5} 推出的四点顶点，那么 \eqref{eq:12.114} 第一行中的因子 $(-i\lambda)$ 就换为
\begin{equation}
  (-2i\lambda)\bigl( \delta^{ij}\delta^{kl} + \delta^{ik}\delta^{jl}
    + \delta^{il}\delta^{jk} \bigr)\cdot\delta^{kl}
  = -2i\lambda(N+2)\delta^{ij} .
\end{equation}
于是
\begin{equation}
  \gamma_{\phi^2} = (N+2)\ \frac{\lambda}{16\pi^2} + \mathcal{O}(\lambda^2) .
  \label{eq:13.50}
\end{equation}

接下来，我们考虑 $(4-\epsilon)$ 维中的同一理论。$\beta$ 函数现在变为
\begin{equation}
  \beta = -\epsilon\lambda + (N+8)\ \frac{\lambda^2}{16\pi^2} ,
  \label{eq:13.51}
\end{equation}
因此在
\begin{equation}
  \lambda^{*} = \frac{16\pi^2\epsilon}{N+8} .
  \label{eq:13.52}
\end{equation}
处存在一个威尔逊-费舍尔不动点。在该不动点处，我们得到
\begin{equation}
  \gamma(\lambda^{*}) = \frac{N+2}{4(N+8)^2}\,\epsilon^2 + \cdots,\qquad
  \gamma_{\phi^2}(\lambda^{*}) = \frac{N+2}{N+8}\,\epsilon + \cdots .
  \label{eq:13.53}
\end{equation}
由这两个结果，我们可以推算出整组临界指数到 $\epsilon$ 阶的预言。作为例子，把 \eqref{eq:13.53} 代入 \eqref{eq:13.19}，得到
\begin{equation}
  \nu^{-1} = 2 - \frac{N+2}{N+8}\,\epsilon + \mathcal{O}(\epsilon^2) ,
  \label{eq:13.54}
\end{equation}
正如第~12.5 节末尾所声称的。

在第~12 章的讨论中，我们声称临界指数的预言与实验数据大致相符。然而，通过计算到更高阶，人们可以获得理论与实验之间精确得多的比较。临界指数的 $\epsilon$ 展开目前已做到 $\epsilon^5$ 阶。更令人瞩目的是，$d = 3$ 时临界指数的 $\lambda$ 展开已做到 $\lambda^9$ 阶。通过对该微扰级数求和，可以得到反常维度 $\gamma(\lambda^{*})$ 和 $\gamma_{\phi^2}(\lambda^{*})$ 非常精确的估计，进而得到临界指数的精确预言。

把这些值与临界指数的直接测定值进行比较，结果列于表~\ref{tab:13.1}。标有``QFT''的列给出由三维 $\phi^4$ 微扰论反常维度计算得到的临界指数值。标有``实验''的列列出在各种系统中临界指数的一些实验测定结果。这些系统包括 Xe、CO$_2$ 及其他流体中的液-气临界点，具有液-液相分离的二元流体混合物中的临界点，Cu-Zn 合金 $\beta$-黄铜原子排列中的有序-无序转变，$^4$He 中的超流转变，以及铁磁体（EuO、EuS、Ni）和反铁磁体（RbMnF$_3$）中的有序-无序转变。不同系统中指数实验测定值之间的一致，是对普适性的直接检验。对于具有单一序参量（$N = 1$）的系统，存在一大批以相同临界指数为特征的、种类繁多的物理系统。

\begin{table}[t]
  \centering
  \caption{三维临界指数：QFT 预言、格点估计与实验的比较。}
  \begin{tabular}{cccc}
    \toprule
    & QFT & 格点 & 实验\\
    \midrule
    $N=1$: $\nu$ & 0.630 & 0.630 & 0.63\\
    \quad $\eta$ & 0.033 & 0.031 & 0.03\\
    $N=2$: $\nu$ & 0.670 & 0.670 & 0.67\\
    \quad $\eta$ & 0.038 & 0.035 & 0.04\\
    $N=3$: $\nu$ & 0.704 & 0.704 & 0.71\\
    \quad $\eta$ & 0.041 & 0.040 & 0.04\\
    \bottomrule
  \end{tabular}
  \label{tab:13.1}
\end{table}

标有``格点''的列包含抽象格点统计力学模型中临界指数的估计值。对于这些简化模型，统计力学配分函数可以在高温展开中计算。经过一番努力，这些展开可以做到高阶并重新求和，从而给出非常精确的临界指数值。表~\ref{tab:13.1} 三列之间——场论、格点与实验——的一致，是对临界现象重整化群描述的惊人确认。

\section{四维临界行为}

本节我们研究 $d = 4$ 的情形，此时 $\phi^4$ 理论恰好可重整。这里无质量 $\phi^4$ 理论的跑动耦合常数在大距离处以对数方式趋于零，但在大动量处增大。Landau 理论的量纲分析预言被对数的幂次修正。这一分析最有趣的应用是解决第~11.4 节末尾提出的问题：$\phi^4$ 理论在 $\mu^2 = 0$ 处的有效势是否具有对称性破缺的极小。

我们首先给出 $\mu^2 = 0$ 的 $\phi^4$ 理论中有效势的单圈表达式。由式~\eqref{eq:11.81}，$\overline{MS}$ 表达式为
\begin{equation}
  V_{eff}(\phi_{cl}) = \frac{\lambda}{4}\phi_{cl}^4\ \Bigl[ 1 +
    \frac{\lambda}{16\pi^2}\ \Bigl( 9\Bigl( \log\frac{\phi_{cl}^2}{M^2}
      - \frac{25}{6} \Bigr) \Bigr) \Bigr] ,
  \label{eq:13.55}
\end{equation}
其中我们令 $N = 1$。该表达式在非零的 $\phi_{cl}$ 处有一个表观极小，但这个极小是微扰展开的假象：它出现在对数足够大以至抵消了领头项的地方，此时微扰论已经失效。

为解决这一问题，我们应用重整化群。跑动耦合常数 $\bar{\lambda}(p)$ 满足
\begin{equation}
  \frac{d}{d\log(p/M)}\ \bar{\lambda} = \frac{3\bar{\lambda}^2}{16\pi^2},
  \qquad \bar{\lambda}(M) = \lambda .
\end{equation}
改进的有效势通过把单圈公式中的跑动耦合和跑动场在由 $\phi_{cl}$ 设定的标度处求值而得到：
\begin{equation}
  V_{eff}(\phi_{cl}) = \frac{\bar{\lambda}(\phi_{cl})}{4}\ \bar{\phi}_{cl}^4
    \ \cdots ,
  \label{eq:13.64}
\end{equation}
其中跑动场 $\bar{\phi}_{cl}$ 包含了反常维度。式~\eqref{eq:13.64} 的误差来自把 $\bar{\lambda}$ 确定为 $\lambda$ 的幂级数。因此该误差为 $\lambda^3$ 阶。当 $\phi_{cl}\to 0$ 时 $\lambda\to 0$，因此表示式 \eqref{eq:13.64} 变得越来越精确。这样该公式成功地求出了危险对数的各次幂之和。把 \eqref{eq:13.64} 视为 $\phi_{cl}$ 的函数，其极小位于 $\phi_{cl} = 0$ 处。因此 \eqref{eq:13.55} 的表观对称性破缺极小确实是不完全微扰展开的假象，在更完整的处理中消失。这解决了我们在第~11.4 节提出的问题。我们应当注意，在更复杂的例子中，微扰论单圈阶发现的有效势的表观对称性破缺极小有可能经得起重整化群分析的检验。第~II 部分的期末课题中给出了一个例子。

我们在此论证中所遵循的程序称为微扰论的\emph{重整化群改进}。这一技术同样可以应用于关联函数的计算以及 Feynman 图微扰论的其他预言：把 Callan-Symanzik 方程的解与按耦合常数同阶的直接微扰论计算结果进行比较，选择重整化群解中未确定的函数，使其复现微扰论的结果。这样，人们得到更紧凑的公式，其中诸如式 (13.56) 中的大对数被重求和进跑动耦合常数。这种重求和产生了关联函数对质量标度对数的依赖，这种依赖刻画了具有边缘或可重整微扰的场论。

在 $\phi^4$ 理论的情形中，跑动耦合常数在小动量处趋于零，在大动量处变大。由于改进微扰论中的误差项是 $\lambda$ 的幂次，改进微扰论在小动量处变得精确，但在大动量处失控。这与我们的物理直觉一致：我们预期微扰论只有在跑动耦合常数保持很小时才精确。

在渐近自由的理论中，跑动耦合常数在大动量处变小，我们可以利用重整化群改进的微扰论在大动量处找到关联函数的精确表达式。在第~17 章和第~18 章中，我们将把这一思想作为分析强相互作用短距离行为的主要工具。

\section{非线性西格马模型}

为了完成对标量场论的研究，我们将讨论一种标量场的非线性理论，其结构与 $\phi^4$ 理论非常不同。这一理论称为\emph{非线性西格马模型}，最初是作为自发对称性破缺的另一种描述而提出的。它之所以引起我们的兴趣，有三个原因。第一，它提供了一个渐近自由理论的简单显式例子。第二，它将给我们第二种维度展开，用它可以研究威尔逊-费舍尔不动点。这样我们就能看到，在远低于 4 的维度 $d$ 中，威尔逊-费舍尔不动点在拉格朗日量空间中走向何处。最后，我们将证明非线性西格马模型在一个不同于标准弱耦合极限的极限中可精确求解。这一解将使我们进一步认识对称性破缺对时空维度的依赖。

\subsection{$d = 2$非线性西格马模型}

我们从二维开始研究。在 $d = 2$ 中，标量场是无量纲的；因此，任何拉格朗日量具有如下形式
\begin{equation}
  \mathcal{L} = \tfrac{1}{2}\sum_i f\bigl( \{\phi^j\} \bigr)\,
    (\partial_\mu\phi^i)^2 ,
  \label{eq:13.65}
\end{equation}
的标量场 $\phi^i$ 理论都具有无量纲耦合，因而是可重整的。由于任何函数 $f$ 都导致一个可重整的理论，这一类标量场论含有无穷多个边缘参数。为限制这些可能的参数，我们必须对理论施加一些对称性。

一个简单的选择是令标量场 $\phi^i$ 构成一个 $N$ 分量单位矢量场 $n^i(x)$，并受约束
\begin{equation}
  \sum_{i=1}^{N} |n^i(x)|^2 = 1 .
  \label{eq:13.66}
\end{equation}
如果我们坚持场论具有 $O(N)$ 对称性，那么 \eqref{eq:13.65} 中的函数 $f$ 就只能依赖于 $n(x)$ 的不变长度，而该长度受 \eqref{eq:13.66} 约束。因此，$f$ 最一般可能的选择是常数。类似地，可以添加到 \eqref{eq:13.65} 上的唯一可能的非导数相互作用 $g(\{n^i\})$ 是常数，而它对 $n$ 的格林函数没有任何影响。有了这些限制，用 $n(x)$ 构建的、具有两个导数且具有 $O(N)$ 对称性的最一般拉格朗日量是
\begin{equation}
  \mathcal{L} = \frac{1}{2g_0^2}\ (\partial_\mu n)^2 .
  \label{eq:13.67}
\end{equation}
该理论具有直接的物理解释。它是这样一个系统的唯象描述：系统的 $O(N)$ 对称性被一个按 $O(N)$ 矢量变换的场的真空期望值自发破缺。考虑，例如，处于自发破缺相的 $N$ 分量 $\phi^4$ 理论中的情形。场 $\phi^i$ 获得一个真空期望值，我们可以把它写成由单位矢量参数化的大小与方向
\begin{equation}
  \langle\phi^i\rangle = \phi_0\,n^i(x) .
  \label{eq:13.68}
\end{equation}
$\phi_0$ 的涨落对应一个质量场，即第~11 章中称为 $\sigma$ 的场。单位矢量 $n(x)$ 方向的涨落对应 $N-1$ 个 Goldstone 玻色子。注意 $n$ 有 $N$ 个分量但受一个约束 \eqref{eq:13.66} 限制，因而含有 $N-1$ 个自由度。形式上，非线性西格马模型是 $\phi^4$ 理论在 $\sigma$ 场质量趋于无穷、同时 $\phi_0$ 保持恒定的极限。

尽管有这种暗示性的联系，我们仍先把非线性西格马模型当作一个独立的量子场论，在自己的立足点上加以分析。求解约束并用 $N-1$ 个 Goldstone 玻色子场 $\pi^k$ 参数化 $n$ 是方便的：
\begin{equation}
  n^i = (\pi^1, \ldots, \pi^{N-1}, \sigma) ,
  \label{eq:13.69}
\end{equation}
其中，按定义，
\begin{equation}
  \sigma = (1-\pi^2)^{1/2} .
  \label{eq:13.70}
\end{equation}
$\pi^k = 0$ 的位形对应沿 $N$ 方向的均匀自发对称性破缺态。表示 \eqref{eq:13.69} 蕴含
\begin{equation}
  (\partial_\mu n)^2 = (\partial_\mu\pi)^2 + \frac{(\pi\cdot\partial_\mu\pi)^2}
    {1-\pi^2} .
  \label{eq:13.71}
\end{equation}
于是拉格朗日量 \eqref{eq:13.67} 取如下形式
\begin{equation}
  \mathcal{L} = \frac{1}{2g_0^2}\ \Bigl[ (\partial_\mu\pi)^2
    + \frac{(\pi\cdot\partial_\mu\pi)^2}{1-\pi^2} \Bigr] .
  \label{eq:13.72}
\end{equation}
注意，$\pi$ 场没有质量项，正如 Goldstone 定理所要求的。

$\pi^k$ 场的微扰论可以通过把拉格朗日量按 $\pi^k$ 的幂次展开而直接读出：
\begin{equation}
  \mathcal{L} = \frac{1}{2g_0^2}(\partial_\mu\pi)^2
    + \frac{1}{2g_0^2}(\pi\cdot\partial_\mu\pi)^2 + \cdots .
  \label{eq:13.73}
\end{equation}
这导致图~\ref{fig:13.1} 所示的费曼规则，外加所有具有偶数个 $\pi^k$ 场的额外顶点。由于拉格朗日量 \eqref{eq:13.67} 是可以用这些场构建的、系数无量纲的最一般 $O(N)$ 对称拉格朗日量，该理论可以通过耦合常数 $g$ 的重整化以及对场 $\pi^k$ 和 $\sigma$ 的 $O(N)$ 对称重新标度而变得有限。在重整化微扰论中，每一个可能的 $2n$-$\pi$ 顶点都有发散和抵消项；然而，这些抵消项都由裸拉格朗日量保持 $O(N)$ 对称性这一基本要求联系起来。

\begin{figure}[t]
  \centering
  \includegraphics[width=0.6\textwidth]{fig13.1.png}
  \caption{非线性西格马模型的费曼规则。}
  \label{fig:13.1}
\end{figure}

我们现在计算该理论的 Callan-Symanzik 函数。由于理论是可重整的，其格林函数对某些函数 $\beta$、$\gamma$ 满足 Callan-Symanzik 方程。明确地，
\begin{equation}
  \Bigl[ M\frac{\partial}{\partial M} + \beta(g)\frac{\partial}{\partial g}
    + n\gamma(g) \Bigr]\ G^{(n)} = 0 ,
  \label{eq:13.74}
\end{equation}
其中 $G^{(n)}$ 是 $n$ 个场 $\pi^k$ 或 $\sigma$ 的格林函数。为了在微扰论领头阶确定 $\beta$ 和 $\gamma$ 函数，我们计算两个简单格林函数到单圈阶，然后看看要使 Callan-Symanzik 方程得到满足，需要什么样的形式。

我们考虑的第一个格林函数是
\begin{equation}
  G^{(1)} = \langle\sigma(x)\rangle .
  \label{eq:13.75}
\end{equation}
展开定义 \eqref{eq:13.70}，得到
\begin{equation}
  \langle\sigma(0)\rangle = \Bigl\langle 1 - \frac{1}{2}\pi^2(0) + \cdots
    \Bigr\rangle = 1 - \frac{1}{2}\langle\pi^2(0)\rangle + \cdots .
  \label{eq:13.76}
\end{equation}
为评估该公式，我们利用图~\ref{fig:13.1} 的传播子计算
\begin{equation}
  \langle\pi^k(0)\pi^\ell(0)\rangle = (N-1)\,\delta^{k\ell}\ \int
    \frac{d^dk}{(2\pi)^d}\ \frac{g_0^2}{k^2+\mu^2} .
  \label{eq:13.77}
\end{equation}
我们加了一个小质量 $\mu$ 作为红外截断。于是
\begin{equation}
  \langle\pi^k(0)\pi^\ell(0)\rangle = (N-1)\ \frac{g_0^2}{(4\pi)^{d/2}}\
    \Gamma(d/2-1)\,\mu^{2-d}\,\delta^{k\ell} .
  \label{eq:13.78}
\end{equation}
把这一结果用于我们 $\langle\sigma\rangle$ 的表达式，然后在动量标度 $M$ 处相减，得到
\begin{equation}
  \langle\sigma(0)\rangle = 1 - \frac{N-1}{2}\ \frac{g^2}{4\pi}\
    \log\frac{M}{\mu} + \cdots .
  \label{eq:13.79}
\end{equation}
只有当
\begin{equation}
  \gamma(g) = \frac{N-1}{4\pi}\,g^2 + \mathcal{O}(g^4) .
  \label{eq:13.80}
\end{equation}
时，该表达式才在 $g$ 阶满足 Callan-Symanzik 方程。

接下来，考虑 $\pi^k$ 两点函数，
\begin{equation}
  \langle\pi^k(p)\pi^\ell(-p)\rangle = \frac{i\delta^{k\ell}}{p^2}
    + \frac{i\delta^{k\ell}}{p^2}\ \Pi^{k\ell}(p)\ \frac{i}{p^2} + \cdots ,
  \label{eq:13.81}
\end{equation}
其中单圈自能 $\Pi^{k\ell}(p)$ 由图~\ref{fig:13.1} 中的费曼规则计算。在评估 $\Pi^{k\ell}$ 时，我们再次遇到积分 \eqref{eq:13.77}，还有积分
\begin{equation}
  \int \frac{d^dk}{(2\pi)^d}\ \frac{i g_0^2\,k^2}{(k^2+\mu^2)^2}
  = \frac{i g_0^2\,d}{(4\pi)^{d/2}}\ \Gamma(2-d/2)\,\mu^{d-2}\ \cdots .
  \label{eq:13.82}
\end{equation}
该公式在 $d = 0$ 处没有极点，并且对于 $d > 0$ 它正比于 $\mu^2$ 的正幂次；因此我们可以令这一收缩为零。于是
\begin{equation}
  \Pi^{k\ell}(p) = \cdots .
  \label{eq:13.83}
\end{equation}
如上在 $M$ 处相减并取极限 $d\to 2$，得到
\begin{equation}
  \langle\pi^k(p)\pi^\ell(-p)\rangle = \frac{i\delta^{k\ell}}{p^2}\ \Bigl[
    1 - \frac{N-1}{4\pi}\,g^2\ \log\frac{M^2}{-p^2} \Bigr] .
  \label{eq:13.84}
\end{equation}
对该结果应用 Callan-Symanzik 方程给出
\begin{equation}
  \Bigl[ M\frac{\partial}{\partial M} + \beta(g)\frac{\partial}{\partial g}
    + 2\gamma(g) \Bigr]\ \langle\pi^k(p)\pi^\ell(-p)\rangle = 0 ,
\end{equation}
由之可以确定 $\beta(g)$：
\begin{equation}
  \beta(g) = -\frac{N-2}{4\pi}\,g^3 + \mathcal{O}(g^5) .
  \label{eq:13.86}
\end{equation}
对于 $N > 2$，$\beta$ 函数的符号为负。这是\emph{渐近自由}的特征：跑动耦合常数在大动量处减小。二维非线性西格马模型是我们遇到的第一个渐近自由理论的例子。

$\beta$ 的负号意味着，在大距离处耦合常数 $g^2$ 增大，微扰展开失效。这也就是说非线性西格马模型中的有序性在长距离处被破坏，这与梅明-瓦格纳定理一致。为了更完整地理解大距离行为，我们可以对 $\beta$ 函数积分：
\begin{equation}
  \bar{g}^2(p) = \frac{16\pi^2}{N-2}\ \frac{1}{\log(p/\Lambda)} ,
  \label{eq:13.87}
\end{equation}
其中 $\Lambda$ 是一个动力学产生的标度。耦合常数在 $p = \Lambda$ 处发散，标志着微扰论的失效以及长距离处有序相的消失。

\subsection{$2 < d < 4$非线性西格马模型}

我们现在把这一分析的结果推广到维度 $d > 2$。对于一般的 $d$，我们继续这样定义非线性西格马模型的作用量：
\begin{equation}
  \int d^dx\ \mathcal{L} = \int d^dx\ \frac{1}{2g_0^2}\ (\partial_\mu n)^2 ,
  \label{eq:13.102}
\end{equation}
其中 $n$ 仍然是无量纲的，因为它满足约束 $|n|^2 = 1$。因此 $g$ 具有量纲 $(\mathrm{mass})^{(2-d)/2}$。我们通过写
\begin{equation}
  T = g^2 M^{d-2} ,
  \label{eq:13.103}
\end{equation}
来定义一个无量纲耦合，正如我们在式 (12.122) 中所做的。如果把 \eqref{eq:13.102} 看作配分函数的 Boltzmann 权重，那么 $T$ 是一个正比于温度的无量纲变量。

由 \eqref{eq:13.103}，我们可以仿照式~\eqref{eq:12.131} 求出 $d$ 维中 $T$ 的 $\beta$ 函数：
\begin{equation}
  \beta(T) = (d-2)T + 2g\beta^{(2)}(g) ,
  \label{eq:13.104}
\end{equation}
其中第二项中的因子 $2g$ 来自定义 $T \sim g^2$。由于 $n$ 是无量纲的，当用无量纲耦合表示时，$\gamma$ 函数与二维结果相同。因此，在 $d = 2 + \epsilon$ 中，
\begin{equation}
  \beta(T) = +\epsilon T - (N-2)\ \frac{T^2}{4\pi};\qquad
  \gamma(T) = (N-1)\ \frac{T}{4\pi} .
  \label{eq:13.105}
\end{equation}
注意，$T$ 的 $\beta$ 函数有一个非平凡的零点，该零点随 $\epsilon\to 0$ 而趋近 $T = 0$。这个零点位于
\begin{equation}
  T^{*} = \frac{4\pi\epsilon}{N-2} .
  \label{eq:13.106}
\end{equation}
$\beta$ 函数的形式示意于图~\ref{fig:13.3}。与第~12.5 节讨论的 $d = 4-\epsilon$ 中的威尔逊-费舍尔零点相反，这是一个\emph{紫外稳定}的不动点。流向红外的流从该不动点出发。由于 $T$ 正比于相应统计系统的温度，$t\to 0$ 是完全有序态，而 $t\to 1$ 是完全无序态。这与伴随 Polyakov 推导 $\beta$ 函数的直觉一致。不动点 $T^{*}$ 对应临界温度。因此，临界温度随 $d\to 2$ 而趋于零，这与梅明-瓦格纳定理一致。

\begin{figure}[t]
  \centering
  \includegraphics[width=0.45\textwidth]{fig13.3.png}
  \caption{非线性西格马模型中 $d > 2$ 时 $\beta$ 函数的形式。}
  \label{fig:13.3}
\end{figure}

我们现在可以在 $\epsilon = d-2$ 的展开中计算非线性西格马模型的临界指数。指数 $\eta$ 直接由
\begin{equation}
  \eta = \gamma(T^{*}) = (N-1)\ \frac{\epsilon}{N-2} .
  \label{eq:13.107}
\end{equation}
给出。为求出第二个指数 $\nu$，我们需要识别对应于 $T \neq T_c$ 时偏离不动点的重整化群流的相关微扰。这正是 $T$ 对 $T^{*}$ 的偏离：
\begin{equation}
  p_T = T - T^{*} .
  \label{eq:13.108}
\end{equation}
由跑动耦合常数的重整化群方程，我们发现跑动的 $p_T$ 满足
\begin{equation}
  \frac{d}{d\log(p/M)}\ p_T = \Bigl( -\epsilon + 2\frac{N-2}{4\pi}T^{*}
    \Bigr)\,p_T .
  \label{eq:13.109}
\end{equation}
方括号中的量为负。如同式~\eqref{eq:12.134} 和 \eqref{eq:12.137}，我们可以把这一量与 $(-1/\nu)$ 等同：在动量 $p \ll M$ 处，
\begin{equation}
  p_T(p) \sim p^{1/\nu} ,
  \label{eq:13.110}
\end{equation}
因此在动量等于 $\xi \sim (T-T^{*})^{-\nu}$ 的倒数处 $p(p)$ 变为 1 量级，正如所要求的。利用 \eqref{eq:13.105} 中 $\beta$ 函数的显式形式，得到
\begin{equation}
  \nu^{-1} = \epsilon ,
  \label{eq:13.111}
\end{equation}
到 $\epsilon$ 的这一阶与 $N$ 无关。（当然，这些结果只适用于 $N > 2$。）热力学临界指数可以用第~13.1 节推导的与模型无关的关系，由 \eqref{eq:13.107} 和 \eqref{eq:13.111} 求出。当这里得到的 $\nu$ 和 $\eta$ 值外推到 $d = 3$（即 $\epsilon = 1$）时，与实验的符合并不出色，但这些结果至少暗示，我们这里找到的不动点可能正是威尔逊-费舍尔不动点向二维附近的延拓。

\subsection{大 $N$ 精确解}

通过用另一种方法处理非线性西格马模型，可以获得对这一不动点性质的进一步认识。由于非线性西格马模型依赖于参数 $N$——单位矢量的分量数——很自然地要问这一模型在 $N\to\infty$ 时表现如何。我们现在证明，如果在保持 $g^2N$ 固定的情况下取这一极限，我们可以得到该模型具有非平凡行为的精确解。

如果我们在欧几里得空间中工作，把拉格朗日量看作自旋系统的 Boltzmann 权重，那么导致这一解的运算最为清晰。于是我们必须计算泛函积分
\begin{equation}
  Z = \int \mathcal{D}n\ \exp\Bigl[ -\int d^dx\ \frac{1}{2g_0^2}
    (\partial_\mu n)^2 \Bigr]\ \prod_x \delta\bigl( n^2(x) - 1 \bigr) .
  \label{eq:13.112}
\end{equation}
这里 $g_0$ 是耦合常数的裸值，而每点一个的 $\delta$ 函数之积实施约束。引入 $\delta$ 函数的积分表示；这需要对拉格朗日乘子变量 $\alpha(x)$ 作第二个泛函积分：
\begin{equation}
  \prod_x \delta\bigl( n^2(x) - 1 \bigr)
  = \int \mathcal{D}\alpha\ \exp\Bigl[ \frac{i}{2g_0^2}\int d^dx\
    \alpha(x)\bigl( n^2(x) - 1 \bigr) \Bigr] .
  \label{eq:13.113}
\end{equation}
现在变量 $n$ 不受约束，并且只以二次形式出现在指数中。因此，我们可以对 $n$ 积分，得到
\begin{equation}
  \begin{aligned}
    Z &= \int \mathcal{D}\alpha\ \bigl( \det[-\partial^2 + i\alpha(x)]
      \bigr)^{-N/2}\ \exp\Bigl[ \frac{i}{2g_0^2}\int d^dx\ \alpha(x) \Bigr]\\
    &= \int \mathcal{D}\alpha\ \exp\Bigl[ -\frac{N}{2}\mathrm{tr}\log
      (-\partial^2 + i\alpha) + \frac{i}{2g_0^2}\int d^dx\ \alpha \Bigr] .
  \end{aligned}
  \label{eq:13.114}
\end{equation}
由于我们在保持 $g^2N$ 固定的情况下取极限 $N\to\infty$，指数中的两项都为 $N$ 量级。因此用最陡下降法评估该积分是有意义的。这意味着用使指数极小的函数 $\alpha(x)$ 的值主导积分。为确定这一位形，我们计算指数对 $\alpha(x)$ 的泛函导数。这给出变分方程
\begin{equation}
  \Bigl( \frac{1}{-\partial^2 + i\alpha} \Bigr)(x,x) = \frac{1}{g_0^2N} .
  \label{eq:13.115}
\end{equation}
该方程的左端必须为常数且为实数；因此，我们应寻找 $\alpha(x)$ 为常数且为纯虚数的解。写
\begin{equation}
  i\alpha = m^2 ,
  \label{eq:13.116}
\end{equation}
则 $m^2$ 满足
\begin{equation}
  N\ \int \frac{d^dk}{(2\pi)^d}\ \frac{1}{k^2 + m^2}
  = \frac{1}{g_0^2} .
  \label{eq:13.117}
\end{equation}
我们将首先在 $d = 2$ 中研究该方程。如果我们用动量截断来定义 \eqref{eq:13.117} 中的积分，就可以评估该积分并求出 $m$ 的方程：
\begin{equation}
  \frac{N}{2\pi}\ \log\frac{\Lambda}{m} = \frac{1}{g_0^2} .
  \label{eq:13.118}
\end{equation}
我们可以通过重整化使该方程有限：
\begin{equation}
  \frac{1}{g_0^2} = \frac{1}{g^2} + \frac{N}{2\pi}\log\frac{\Lambda}{M} ,
  \label{eq:13.119}
\end{equation}
该式引入了一个任意的重整化标度 $M$。于是我们可以解出 $m$，得到
\begin{equation}
  m = M\ \exp\Bigl( -\frac{2\pi}{N g^2} \Bigr) .
  \label{eq:13.120}
\end{equation}
质量 $m$ 是 $\pi$ 场动力学产生的质量。它是一个非微扰量：由于 $g^2$ 出现在指数中，它到耦合常数的所有阶都为零。这一质量的存在意味着 Goldstone 玻色子通过非微扰效应获得质量，从而在长距离处破坏有序态。这是梅明-瓦格纳定理在大 $N$ 极限下的精确实现。

在 $2 < d < 4$ 中，同样的分析给出非平凡的不动点和动力学产生的质量，证实了从重整化群得到的图景。对于 $d > 4$，方程 \eqref{eq:13.117} 只有平凡解 $m = 0$，有序性在大距离处得以保持。

我们现在已经获得了 $O(N)$ 对称标量场论作为维度函数的行为的完整图景。对于 $d > 4$，理论由自由场不动点描述，临界行为是 Landau 理论的临界行为。在 $d = 4$ 处，理论是可重整的，临界行为被对数修正。对于 $2 < d < 4$，理论由威尔逊-费舍尔不动点控制，临界指数与 Landau 值相差可计算的修正。对于 $d = 2$，理论由非线性西格马模型描述，它是渐近自由的，在长距离处不表现出有序性。模型行为随 $d$ 的这种演化阐明了前两章的主要观点：量子场论的定性行为不是由基本拉格朗日量决定，而是由重整化群流及其不动点的性质决定。而这些反过来又只依赖于施加于相互流动的拉格朗日量族上的基本对称性。这一结论在最深层次上标志着对称性原理在确定基本物理规律中的重要性。

\problems

\textbf{13.1} \emph{修正标度的指数。} 对于 $4-\epsilon$ 维中的临界现象，消失得最慢的无关贡献与耦合常数 $\lambda$ 对其不动点值的偏离相关。这给出了对第~13.1 节导出的标度律的最重要的非普适修正。通过研究 Callan-Symanzik 方程的解，证明：如果 $\lambda$ 的裸值与 $\lambda^{*}$ 略有不同，则 Gibbs 自由能会得到一个修正
\begin{equation}
  G(M,t) \to G(M,t)\cdot\Bigl( 1 + (\lambda-\lambda^{*})\,
    t^{-\omega}\ \cdots \Bigr) .
\end{equation}
该公式定义了一个新的临界指数 $\omega$，称为\emph{修正标度的指数}。证明
\begin{equation}
  \omega = \nu\,\beta'\bigl( \lambda^{*} \bigr) .
\end{equation}

\textbf{13.2} \emph{指数 $\eta$。} 结合习题~10.3 的结果与适当的重整化方案，证明 $\phi^4$ 理论中 $\gamma(\lambda)$ 的领头项为
\begin{equation}
  \gamma = \frac{\lambda^2}{12(4\pi)^4} .
\end{equation}
将这一结果推广到 $O(N)$ 对称 $\phi^4$ 理论以推导式~\eqref{eq:13.47}。计算 $\eta$ 的领头阶（$\epsilon^2$）贡献。

\textbf{13.3} \emph{$CP^N$ 模型。} 正文中讨论的非线性西格马模型可以看作场为单位球面上坐标的量子理论。一个稍复杂的高对称空间是复射影空间 $CP^N$。该空间可以定义为满足条件
\begin{equation}
  \sum_j |z_j|^2 = 1 ,
\end{equation}
的 $(N+1)$ 维复矢量 $(z_1, \ldots, z_{N+1})$ 的空间，其中相差一个整体相位旋转的点被视为等同，即
\begin{equation}
  (e^{i\alpha}z_1, \ldots, e^{i\alpha}z_{N+1}) \quad
  \text{identified with} \quad (z_1, \ldots, z_{N+1}) .
\end{equation}
在本习题中，我们研究场为该空间上坐标的二维量子场论。
\begin{enumerate}[label=(\alph*)]
  \item 表示 $CP^N$ 上坐标理论的一种方式是写出一个依赖于受约束的场 $z_j(x)$ 的拉格朗日量，该拉格朗日量还具有局域对称性
    \begin{equation}
      z_j(x) \to e^{i\alpha(x)} z_j(x) ,
    \end{equation}
    在每个点 $x$ 处独立。证明如下拉格朗日量具有这一对称性：
    \begin{equation}
      \mathcal{L} = \frac{1}{g^2}\ \Bigl[ |\partial_\mu z_j|^2
        - |z_j^{*}\partial_\mu z_j|^2 \Bigr] .
    \end{equation}
    为证明不变性，你将需要用到对 $z_j$ 的约束及其推论
    \begin{equation}
      z_j^{*}\partial_\mu z_j = -(\partial_\mu z_j^{*})z_j .
    \end{equation}
    证明 $N = 3$ 情形的非线性西格马模型可以通过替换
    \begin{equation}
      n^i = z^{*}_j(\sigma^i)^j_{\ k}\,z^k ,
    \end{equation}
    转换为 $N = 1$ 情形的 $CP^N$ 模型，其中 $\sigma^i$ 是 Pauli sigma 矩阵。
  \item 为把拉格朗日量写成更简单的形式，引入一个实施约束的标量拉格朗日乘子 $\lambda$，以及一个表示局域对称性的矢量拉格朗日乘子 $A_\mu$。更具体地，证明 $CP^N$ 模型的拉格朗日量可以由拉格朗日量
    \begin{equation}
      \mathcal{L} = \frac{1}{g^2}\ \Bigl[ |D_\mu z_j|^2 - \lambda\bigl(
        |z_j|^2 - 1 \bigr) \Bigr] ,
    \end{equation}
    （其中 $D_\mu = (\partial_\mu + iA_\mu)$）对场 $\lambda$ 和 $A_\mu$ 作泛函积分而得到。
  \item 我们可以通过对场 $z_j$ 积分来求解 $N\to\infty$ 极限下的 $CP^N$ 模型。证明该积分导致表达式
    \begin{equation}
      Z = \int \mathcal{D}\lambda\,\mathcal{D}A\ \exp\Bigl[ -iN\,
        \mathrm{tr}\log\bigl( -D^2 - \lambda \bigr) + \frac{iN}{g^2}\int d^2x\
          \lambda \Bigr] ,
    \end{equation}
    其中我们只保留了 $N\to\infty$、$g^2N$ 固定时的领头项。用类似于我们用于非线性西格马模型的方法，考察指数对 $\lambda$ 和 $A_\mu$ 极小化的条件。证明这些条件在 $A_\mu = 0$ 和 $\lambda = m^2 > 0$ 处有解。证明：如果 $g^2$ 在标度 $M$ 处重整化，$m$ 可以写为
    \begin{equation}
      m = M\ \exp\Bigl( -\frac{2\pi}{Ng^2} \Bigr) .
    \end{equation}
  \item 现在把指数在 $A_\mu = 0$ 附近展开。证明该展开中第一个非平凡的项正比于有质量标量场的真空极化。用维数正规化评估该表达式，并证明它为 $A_\mu$ 给出一个标准的动能项。这样，我们最初出发的那个奇特非线性场论最终转化为 $(N+1)$ 个有质量标量场与一个无质量光子相互作用的理论。
\end{enumerate}
